% Source: Weinberg GR, physical PDF pp. 230--231
%         (printed pp. 207--208).
% Coverage: Section 8.8, The Schwarzschild Singularity; equations
%           (8.8.1)--(8.8.3).
% Figures/tables/footnotes: one optional-reading footnote; reference marker 20.
% Status: source-reviewed and compile-clean.
% Uncertainties: Eq. (8.8.3) is converted from the source's time-positive
%                local coordinate ordering to the target (-+++) convention;
%                see MODERNIZATION CHECK below.

\subsection{The Schwarzschild Singularity\optionalreading}\label{sec:8.8}
\begingroup
\renewcommand{\thefootnote}{*}
\begin{NoHyper}
\footnotetext{This section lies somewhat out of the book's main line of
development, and may be omitted in a first reading.}
\end{NoHyper}
\endgroup

The reader will probably have noticed that the Schwarzschild solution
\eqref{eq:8.2.12} becomes singular at \(r=2MG\). This radius corresponds
to \(\rho=MG/2\) and \(R=MG\), so we see that this singularity also occurs
when the metric is expressed in its isotropic form \eqref{eq:8.2.14} or in
its harmonic form \eqref{eq:8.2.15}. The radius \(2GM\) at which the
singularity occurs in standard coordinates is called the
\emph{Schwarzschild radius} of the mass \(M\).

It should immediately be stressed that there is no Schwarzschild
singularity in the gravitational field of any known object in the universe.
The Schwarzschild singularity appeared in the solution of Einstein's vacuum
equations \(R_{\mu\nu}=0\), and is therefore irrelevant if the radius
\(2GM\) lies within the massive body, where we must use the full Einstein
equation \eqref{eq:7.1.13}. For the sun the Schwarzschild radius
\(2GM_\odot\) is 2.95 km, deep in the solar interior, and we shall see in
Chapter~11 that the solution of the full Einstein equation inside a stable
star exhibits no Schwarzschild singularity, or any other singularity. For
a proton the Schwarzschild radius is \(10^{-50}\,\mathrm{cm}\), and this is
37 orders of magnitude smaller than the characteristic proton radius of
about \(10^{-13}\,\mathrm{cm}\)! In Chapter~11 we discuss the possibility
that a very massive body might collapse to a radius smaller than its
Schwarzschild radius, but with this one hypothetical exception the
Schwarzschild singularity does not seem to have much relevance to the real
world.

It is nonetheless instructive to imagine a body so small and massive that
the radius \(2GM\) lies outside it, in empty space. The Schwarzschild
solution then holds down to this radius and actually displays a singularity.
But is this singularity real? We can readily calculate the four
nonvanishing curvature invariants described in Section~6.7, and find that
they are all perfectly well behaved at the Schwarzschild radius although
they do become singular at the origin. This suggests that the apparent
Schwarzschild singularity may be only an artifact of the coordinate systems
we have used. (If any one of the curvature invariants had been singular at
the Schwarzschild radius, then this singularity would of course have been
present in all coordinate systems.) Only a few years ago a coordinate
system was found that allows us to avoid talking about a Schwarzschild
singularity, if we are willing to allow the world an unusual
topology.\hyperref[ch8-ref-20]{\textsuperscript{20}} To exhibit this
reinterpretation of the Schwarzschild singularity, we introduce a new set
of coordinates \(r'\), \(\theta\), \(\varphi\), \(t'\), defined by
\begin{equation}
  r'^2-t'^2
  \equiv
  T^2\left(\frac{r}{2GM}-1\right)
  \exp\left(\frac{r}{2GM}\right),
  \tag{8.8.1}\label{eq:8.8.1}
\end{equation}
\begin{equation}
  -\frac{2r't'}{r'^2+t'^2}
  \equiv
  \tanh\left(\frac{t}{2MG}\right),
  \tag{8.8.2}\label{eq:8.8.2}
\end{equation}
where \(T\) is an arbitrary constant. The Schwarzschild solution
\eqref{eq:8.2.12} then becomes
% MODERNIZATION CHECK: in the local (r',theta,phi,t') discussion the source
% prints d\tau^2=C(dt'^2-dr'^2)-r^2d\Omega^2.  Translating to the target
% convention ds^2=-d\tau^2 and writing the time differential first gives
% the (-+++) form below.
\begin{equation}
  \dd s^2
  =
  \left(\frac{32G^3M^3}{rT^2}\right)
  \exp\left(-\frac{r}{2GM}\right)
  (-\dd t'^2+\dd r'^2)
  +r^2\dd\theta^2+r^2\sin^2\theta\,\dd\varphi^2,
  \tag{8.8.3}\label{eq:8.8.3}
\end{equation}
where \(r\) is now to be understood as a function of \(r'^2-t'^2\)
defined by Eq.~\eqref{eq:8.8.1}. The metric is nonsingular as long as
\(r^2\) is well defined and positive definite, that is, as long as
\[
  r'^2>t'^2-T^2.
\]
Hence, during the time interval \(0<t'<T\), the metric is a perfectly
smooth finite function of \(r'\) for all real \(r'\). Indeed, even
\(g_{\theta\theta}\) and \(g_{\varphi\varphi}\) do not vanish when
\(r'=0\), so that when we approach the origin \(r'=0\) there is nothing to
keep us from continuing right through to negative \(r'\)! The space
described by Eq.~\eqref{eq:8.8.3} is therefore singularity free but
consists of two identical sheets \(r'>0\) and \(r'<0\), joined in a smooth
way by a branch point at \(r'=0\). When \(t'\) reaches the time \(T\) the
two sheets detach from each other, and thereafter the metric has a real
singularity at
\[
  r'=\pm\sqrt{t'^2-T^2},
\]
that is, at \(r=0\). However, even so, the metric has no singularity at the
radius \(r'=t'\) that corresponds to the Schwarzschild radius \(r=2GM\).

To repeat, this discussion of the Schwarzschild singularity does not apply
to any gravitational field actually known to exist anywhere in the
universe. Indeed, it does not even apply to gravitational collapse (see
Section~11.9), because for \(t'<T\) space is empty for all \(r'\). However,
like Aesop's fables, it is useful because it points to a moral, that what
appears in one coordinate system to be a singularity may in another
coordinate system have quite a different interpretation.
